\documentclass[11pt,a4paper]{article}
\usepackage[T1]{fontenc}
\usepackage{lmodern,amsmath,amssymb,amsthm,booktabs,microtype}
\usepackage[margin=24mm]{geometry}
\usepackage[colorlinks=true,urlcolor=blue,linkcolor=blue]{hyperref}
\newcommand{\Js}{\mathcal J_s}
\newcommand{\E}{\mathbb E}
\newcommand{\tr}{\operatorname{tr}}
\newcommand{\av}[1]{\langle #1\rangle}
\newcommand{\W}{\mathcal W}
\newcommand{\JH}{\mathsf J_H}
\title{Reducing the raw cubic singlet\\Conserved polynomials, descendants and cross spectra}
\author{Alan Day \and Chengheng\thanks{AI research assistant (OpenAI).}}
\date{Local research note --- 6 November 2026}
\begin{document}\maketitle
\begin{abstract}
The raw cubic derivative singlet contains several components with the
same disorder. We derive exact finite-size identities that reduce its
cubic energy block to a conserved polynomial and its mixed block to an
energy-weighted quadratic singlet. Their linear counterterms vanish
as $1/p$ in double scaling; the leading energy-block cross correlation
is fixed by the known linear-residual kernel. An exhaustive endpoint
certificate also removes the cubic-residual--quadratic cross spectrum.
The remaining spectral problem concerns the linear residual's cross
spectra with the quadratic and cubic channels.
\end{abstract}
\section{The reduced physical question}
Use standard CAR and the complex Gaussian number-preserving $p$-body
Hamiltonian $H=H_s$, including overlapping index sets. Brackets mean
ensemble average of normalized trace, at $\beta=\mu_{\rm chem}=0$.
Set $L=[H,\cdot]$, $\JH B=\{H,B\}/2$, and $\hbar=1$; time is
conjugate to $H_s$. First take $N,p\to\infty$, with
\[
 \lambda=p^2/N>0,\quad q=e^{-4\lambda},\quad
 \sigma^2=\Js^2/\lambda .
\]
The preceding notes \cite{local} define
\[
 A=\sum_i(Lc_i)^\dagger Lc_i,\quad M_3=\sum_i c_i^\dagger L^3c_i,
 \quad T_3=(M_3+M_3^\dagger)/2,
\]
\[
 (T,S,Z)=\W_3(T_3,\JH A,H^3),\qquad R=T-aS-bZ,
\]
where $\W_j$ is ordinary Gaussian \emph{coupling} chaos, and
\[
 a=-\frac{(1-q)^2}{1+q^2},\qquad
 b=\frac{2\lambda(1-q)(1+q^3)}{(1+q^2)(1+q+q^2)}.
\]
After subtracting an explicit conserved polynomial $P_N(H,Q_f)$,
our exact reduction and its double-scaled limit are
\begin{align}
 M_3-P_N&=R+a\JH Y_N+F+\tfrac32 LY_N-aB_S-bB_Z,\label{raw}\\
 B_S,B_Z&\longrightarrow0\quad\hbox{in ensemble Hilbert--Schmidt norm}.
\end{align}
Here $F=\W_1T_3-\bar t(Q_f)H$ and $Y_N$ is the projected quadratic
singlet. Their definitions and the proof follow. The diagonal kernels
of $F,Y,R$ have weights $q,q^2,q^3$ respectively; they alone do not
determine the spectrum of (\ref{raw}).

\newpage
\section{Exact first-chaos channels}
Work in charge sector $m$, with $H=\sum_\nu H_\nu$ the orthogonal
irreducible $U(N)$ decomposition. Set
\begin{align}
 C&=w\binom mp\binom{N-m+p}p,\qquad C_-=C_{m-1},\\
 \kappa_\nu&=w\binom{m-\nu}p\binom{N-m+p-\nu}p,\qquad
 \av{H_\nu^2}=C\omega_\nu,\quad \sum_\nu\omega_\nu=1.
\end{align}
$w$ is the coupling covariance; it is not $\lambda$. The exact rank
weights and covariance eigenvalues are the embedded-ensemble formulas
\cite{kota}, also used in our earlier channel calculation.
Three Wick contractions give
\begin{equation}
 \W_1H^3=\sum_\nu z_\nu H_\nu,\qquad z_\nu=2C+\kappa_\nu.\label{z}
\end{equation}
For the Jordan block the result is
\begin{align}
 \W_1\JH A&=\sum_\nu s_\nu H_\nu,\qquad s_\nu=CK_0+\kappa_\nu J_\nu,\label{s}\\
 K_0&=2p-m+mC_-/C+p^2/m,\\
 J_\nu&=\frac{p[p(N-m+1)-\nu(N-\nu+1)]}
 {(m-\nu)(N-m-\nu+1)}.
\end{align}
One derivation uses $\mathfrak D B=Q_fB-\sum_i c_i^\dagger Bc_i$
and the exact identities
\[
 A=(2p-\mathfrak D)H^2,\qquad
 T_3=(3p-\mathfrak D)H^3-3\JH A.
\]
The embedded rank-$\nu$ tensor obeys
\[
 \bigl(\mathfrak D[z_\nu(Q_f)H_\nu]\bigr)_m
 =[mz_\nu(m)-(m-p)z_\nu(m-1)]H_{\nu,m}.
\]
Consequently the previously derived $t_\nu$ for $\W_1T_3$ gives
\begin{equation}
 s_\nu=\frac{(3p-m)z_\nu(m)+(m-p)z_\nu(m-1)-t_\nu(m)}3.\label{safe}
\end{equation}
This nonsingular formula also covers the removable boundary in
(\ref{s}). Define deterministic sector averages and centered operators
\begin{align}
 \bar z&=\sum_\nu\omega_\nu z_\nu,& B_Z&=\sum_\nu(z_\nu-\bar z)H_\nu,\\
 \bar s&=\sum_\nu\omega_\nu s_\nu,& B_S&=\sum_\nu(s_\nu-\bar s)H_\nu.
\end{align}
Thus $Z=H^3-\bar zH-B_Z$. At finite size $Z$ need not commute
with $H$: coupling-chaos subtraction is not the same operation as
forming an orthogonal polynomial in $H$. The latter is the
$q$-Hermite construction in \cite{tall}.

\newpage
\section{Counterterms and the first finite-size scale}
At half filling put $s=p-\nu$. Earlier exact rank ratios give
$s\Rightarrow\operatorname{Pois}(\lambda)$ with a uniform factorial
tail, and $\|F\|_2^2\to V=64\Js^6q^2$. We now obtain
\begin{equation}
 \boxed{\left\|pB_Z-\tfrac12F\right\|_2\to0,\qquad
 \left\|pB_S-\tfrac{1-2\lambda}{2}F\right\|_2\to0.}\label{scale}
\end{equation}
Here $\|B\|_2^2=\E\tr(B^\dagger B)$. To see the coefficients,
write $k_\nu=\kappa_\nu/C$. The exact product is
\[
 k_\nu=\prod_{j=0}^{p-1}
 \left(1-\frac{\nu}{m-j}\right)
 \left(1-\frac{\nu}{N-m+p-j}\right).
\]
For each fixed $s$, expanding this product and the rational $J$ gives
\begin{align}
 k_p&\to q,&p(k_{p-s}-k_p)&\to4\lambda qs,\\
 J_p&\to-2\lambda,&p(J_{p-s}-J_p)&\to4\lambda s,\\
 p\frac{s_{p-s}-s_p}{C}&\to4\lambda q(1-2\lambda)s.
\end{align}
The corresponding previously established slope is
$(t_{p-s}-t_p)/C\to8\lambda qs$. Center all three coefficients
under the \emph{finite} rank weights, then pass to the Poisson limit.
This yields (\ref{scale}), including its constants.

For completeness the tail step can be controlled without a rank
cutoff. In the central charge window $|m-N/2|\leq N^{2/3}$,
all denominators in the product and in $J$ are bounded below by
$cN$ for sufficiently large $N$. For real $0\leq\nu\leq p$,
\[
 0\leq k_\nu\leq1,\quad
 |\partial_\nu k_\nu|\leq K_\lambda/p,\quad
 |J_\nu|\leq K_\lambda,\quad
 |\partial_\nu J_\nu|\leq K_\lambda/p.
\]
Hence $p|k_{p-s}-k_p|$ and
$p|k_{p-s}J_{p-s}-k_pJ_p|$ are bounded by $K_\lambda s$.
Together with the earlier bound for the $t$ slope and
$\omega_{p-s}\leq K_\lambda A_\lambda^s/s!$, this proves
weighted $L^2$ convergence and justifies centering. The expansions
are uniform on this charge window. Outside it, the exact nonsingular
coefficients have polynomial growth times the maximal sector variance;
the binomial charge tail is exponentially small. This extends
(\ref{scale}) to the full infinite-temperature trace.

The limit is at fixed positive $\lambda$, followed, if needed, by
$\lambda\to0^+$. No simultaneous $p,\lambda,t$ scaling is asserted.

\newpage
\section{A nonzero cross term and its controlled limit}
The conserved part of $Z$ is $H^3-\bar z(Q_f)H$. Since
$\av{FZ}=0$ by ordinary Gaussian-chaos orthogonality, invariance
of trace under evolution gives the exact all-time identity
\begin{equation}
 \boxed{\av{F(t)Z}=\av{FB_Z}-\av{F(t)B_Z}.}\label{fz}
\end{equation}
It follows that $|\av{F(t)Z}|\leq2\|F\|_2\|B_Z\|_2$, uniformly
in real time. Thus the unscaled cross term vanishes in DS, but its
leading correction is nontrivial:
\begin{equation}
 \boxed{\lim_{\rm DS}2p\av{F(t)Z}=V[1-C_F(t)]},\qquad
 C_F(t)=\frac{\av{F(t)F}_{\rm DS}}V.\label{leading}
\end{equation}
The residual after replacing $2pB_Z$ by $F$ in (\ref{fz}) is
bounded by $2\|F\|_2\|2pB_Z-F\|_2$. The remaining time limit
is the previously proved $r=q$ Mehler kernel. Equation
(\ref{leading}) is therefore a fixed-time statement, obtained
without differentiating a merely pointwise time limit.

The curvature follows independently. For any invariant linear
Gaussian filter $B$, tracial H\"older and Gaussian fourth-moment
bounds imply
\[
 \|LB\|_2^2\leq4\|H\|_4^2\|B\|_4^2
 \leq12C\|B\|_2^2.
\]
Applying this to the error in (\ref{scale}) proves
\begin{equation}
 2p\av{F L^2Z}\longrightarrow-\av{F L^2F}_{\rm DS}
 =-\frac{128\Js^8q^2(1-q)}\lambda.\label{m2}
\end{equation}
No independence between $B$ and $H$ is used.

An exact finite example demonstrates why the cross terms had to be
computed. With $N=5,p=2,m=3,w=1$, let
$B_u(C_3)=\av{FH^u C_3 H^{2-u}}$ for $C_3=T,S,Z$. Full
polarized CAR tensors and the six allowed Wick contractions give
\begin{center}
\begin{tabular}{lrrrr}\toprule
Block & $B_0$ & $B_1$ & $B_2$ & $\av{F L^2 C_3}$\\\midrule
$T$ &4590&5400&4590&$-1620$\\
$S$ &$-6570$&$-5760$&$-6570$&$-1620$\\
$Z$ &$-4680$&$-3600$&$-4680$&$-2160$\\\bottomrule
\end{tabular}
\end{center}
Here $\av{F^2}=180$ and $\av{FL^2F}=3240$. These finite-size
values are not obtained by substituting $p=2$ into a DS asymptotic.

\newpage
\section{The quadratic descendant and the raw operator}
Define the exact sector projection
\begin{align}
 a_0&=\E A=(2p-m)C+mC_-,\qquad O=A-a_0\mathbf1,\\
 Y_N&=O-c_N(H^2-C\mathbf1),\qquad
 c_N=\frac{\av{O(H^2-C)}}{\av{(H^2-C)^2}}.
\end{align}
All scalar coefficients act as functions of $Q_f$ on full Fock
space. In sectors with zero denominator set $c_N=0$; these
degenerate sectors do not contribute to the central DS limit.
For nondegenerate sectors, $\av{H\JH A}=C\bar s$ and the three
fourth-order Wick pairings give
\begin{align}
 c_N&=\frac{\bar s-a_0}{C+\bar\kappa}
 =\frac{p^2/m+\sum_\nu\omega_\nu(\kappa_\nu/C)J_\nu}
 {1+\sum_\nu\omega_\nu\kappa_\nu/C}\longrightarrow c,\\
 c&=\frac{2\lambda(1-q)}{1+q},\qquad
 \bar\kappa=\bar z-2C.\label{c}
\end{align}
Because $\av{H^2Y_N}=0$, the first chaos of its Jordan descendant
has no scalar $H$ component:
\begin{align}
 G_N:=\W_1\JH Y_N&=B_S-c_NB_Z,\\
 S&=c_N Z+\JH Y_N-G_N,\label{sreduce}\\
 pG_N&\longrightarrow\tfrac12(1-2\lambda-c)F
 \quad\hbox{in }L^2.
\end{align}
Indeed, apply $\W_3$ to
$\JH Y_N=\JH A-c_NH^3+(c_NC-a_0)H$, and use
$\bar s-a_0=c_N(\bar z-C)$. This proves (\ref{sreduce}) exactly.

Using $T_3=R+aS+bZ+\bar tH+F$ and
$M_3=T_3+\tfrac32 LA$, set
\begin{equation}
 P_N=(ac_N+b)H^3+[\bar t-(ac_N+b)\bar z]H.\label{poly}
\end{equation}
Then substitution yields (\ref{raw}) with no omitted finite-size
terms; also $LA=LY_N$. Thus the independent energy-resolved
unknowns can be organized in the basis $(F,Y,R)$, rather than
$(F,S,Z,R)$. In DS,
\[
 Y=\sqrt{W_d}X,\qquad
 W_d=4\Js^4\left[\lambda^{-1}+\frac{4q}{1+q}\right],
\]
where $X$ is our earlier normalized quadratic singlet. In an
energy basis the $Y$ term in (\ref{raw}) is weighted by
\[
 a\frac{E+E'}2+\frac32(E-E').
\]
The raw $M_3$ is not Hermitian: the spectral Gram matrix must use
the appropriate adjoint. This reduction specifies the observable
before its cross spectra are combined.

\newpage
\section{The cubic residual decouples from the quadratic spectrum}
The earlier exhaustive endpoint reduction for two cubic blocks can
be reused with different weights. Put
\[
 D=(1+q^2)(1+q+q^2),\quad B=2(1-q)(1+q^3),
 \quad(P_1,P_2,P_3)=(-q(1+q+q^2),-q(1+q+q^2),D).
\]
For a pairing $\pi$, let $\mathrm{cr}$ count crossings, let
$\alpha_i,\beta_j=2-4\mathbf1_{\rm crossing}$ for the marked
chord pairs $12,23,13$ at the respective blocks, and let
$\delta_{ij}=4$ for a nested double bridge and zero otherwise.
Up to common nonzero factors, the $R$--$S$ orbit has coefficients
\begin{align}
 [\lambda]:&\quad\sum_\pi q^{\mathrm{cr}}
                   \sum_i P_i(\delta_{i1}+\delta_{i2}),\\
 [\lambda^2]:&\quad\sum_\pi q^{\mathrm{cr}}
              (\textstyle\sum_iP_i\alpha_i-B)(\beta_1+\beta_2),
\end{align}
and $R$--$Z$ has coefficient
$\sum_\pi q^{\mathrm{cr}}(\sum_iP_i\alpha_i-B)$.
Every one is identically zero in every endpoint template,
including the fully bridged template.

To explain the all-order scope, delete spectator chords having both
endpoints outside the two consecutive cubic blocks. Permuting slots
within either block leaves its crossing count with every spectator
unchanged. The remaining template is specified by bridge count
$k=0,1,2,3$, the color order of $6-2k$ external endpoints, and the
split between the two arcs. There are $140+30+6+1=177$ templates
and $6234$ pairings. Exact polynomial evaluation of these coefficients
therefore proves, for every fixed $u,v\geq0$,
\begin{equation}
 \av{RH^uSH^v}_{\rm DS}=\av{RH^uZH^v}_{\rm DS}=0.\label{rs}
\end{equation}
The local minor factors are separately checked against the existing
determinant implementation in all $37404$ marked cases. This shares
the earlier template generator, rather than constituting a second
independent enumeration proof.

There is a useful consequence for $Y$. For Hermitian insertions
define the matrix of two-energy measures by
$d\nu_{AB}=\E\tr[P(dE)A P(dE')B]$. It is positive as a matrix
measure and satisfies spectral Cauchy--Schwarz. The diagonal $R,Y$
limits are the absolutely continuous Mehler measures with weights
$q^3,q^2$. Their moment bounds give tightness and uniform
integrability for cross measures, supported in the same limiting
compact energy square. The $L^2$-small linear $G_N$ makes no fixed
polynomial moment contribution (Gaussian moment bounds and H\"older).
Equations (\ref{sreduce}) and (\ref{rs}) then imply
$(E+E')d\nu_{RY}=0$. Spectral Cauchy--Schwarz makes $\nu_{RY}$
absolutely continuous; it cannot be supported on the line $E+E'=0$.
Thus $\boxed{d\nu_{RY}=0}$. This argument applies to every subsequential
limit, establishing the full vanishing cross-spectrum limit.

\newpage
\section{Evidence, remaining calculation and sources}
The remaining unknown energy-resolved cross measures are
$\nu_{FR}$ and $\nu_{FY}$. In particular, an ordinary time correlator
that vanishes by parity would not dispose of the energy-sum-weighted
$F$--$\JH Y$ term. The next calculation should insert the centered
rank filter $F$ into the existing marked-chord contractions, keeping
the first correction that survives its rank centering. The present
result does not yet justify summing three diagonal probe spectra.

The new numerical records in
\texttt{results/susskind-cubic-cross} contain:
\begin{itemize}
\item Three finite CAR systems for $F$ against $T,S,Z$ with two external
Hamiltonians, and 32 first-chaos Jordan-channel comparisons from
four $(N,p)$ systems. The finite tensor computation uses bounded
Gaussian-integer arrays and rational final normalization.
\item Fifteen full-rank sums, $\lambda=0.1,0.25,0.5$ and
$p=16,32,64,128,256$, repeated at 80 and 120 decimal digits.
Their stored scaled-counterterm squared errors agree within
$10^{-60}$; this is a precision check, not interval certification.
\item The exact 177-template cross certificate and 37404 local
determinant-minor checks described above.
\item An independent rank-bilinear calculation reproduces the three
finite $F$--$Z$ curvatures, including $-2160$. Twelve exact rational
matrix identities check the Jordan, raw-operator and commutator
reductions in three deterministic systems. An initial checker used
structural equality on unexpanded expressions and stopped; expanding
the exact residuals resolves it. The original checker is retained.
\end{itemize}
For example, at $\lambda=0.1,p=256$ the squared errors in the
$Z,S,G$ versions of (\ref{scale}), divided by $\|F\|_2^2$, are
respectively $6.54\times10^{-6}$, $8.04\times10^{-6}$ and
$7.48\times10^{-6}$. At $\lambda=0.5,p=256$ the $S$ error is
$1.94\times10^{-5}$. All ranks are summed; finite-size deviations
are retained, and these numbers are not errors on a bulk prediction.

The known covariance and $q$-Hermite constructions below are used as
inputs, not claimed as new results. The new work in this note is the
explicit reduction for these same-coupling derivative insertions and
its stated cross-spectrum consequences. It remains a model calculation
within the Susskind singlet program: bulk locality, masses and the
finite-gravity interaction dictionary require further observables.

\begin{thebibliography}{9}
\bibitem{local} A. Day and Chengheng, local research notes (2026):
\emph{An independent cubic derivative singlet};
\emph{The missing linear component of the cubic singlet};
\emph{Dynamics of the linear cubic residual};
\emph{The linear cubic residual has a full $q$ kernel};
and the quadratic/cubic full-kernel notes. Source directories:
\texttt{manuscript/susskind-cubic-*},
\texttt{susskind-linear-kernel}, \texttt{susskind-singlet-kernel}.
\bibitem{kota} V. K. B. Kota,
\emph{Bivariate moments of the two-point correlation function for
embedded Gaussian unitary ensemble with $k$-body interactions},
\href{https://arxiv.org/html/2208.11312}{arXiv:2208.11312v2},
Sec. IV, Eqs. (28)--(31).
\bibitem{tall} J. Tall and S. Tomsovic,
\emph{Double-scaled bosonic and fermionic embedded ensembles,
complex SYK, and the dual Hilbert space},
\href{https://arxiv.org/html/2604.14522v2}{arXiv:2604.14522v2},
Sec. 3.1, Eqs. (3.1)--(3.14), and Sec. 4.1.
\end{thebibliography}
\end{document}
